% ============================================================================
% SAMPLE TEXT — /thesis-review v3, section 5.1 `fi-why` (2026-08-21)
% Full-section sample on Matteo's explicit request: five-move architecture
% (structure -> problem+posthoc -> grounding-as-forcing -> payoff -> hand-off)
% then the asymmetry part (his surviving text with agreed repairs).
% SAMPLE ONLY: never paste into chapters/*.tex — adapt in your own words.
% Chapter lead-in deliberately absent (parked; to be redone later).
% TODO-cites: melchiorre2022protomf + the sleep-stages paper (Niknazar &
% Mednick original) have no bibliography.bib entries yet.
% LLM-usage-disclosure artifact.
% ============================================================================

% [lead-in placeholder — parked]

% --- PART ONE: the mechanism ---

% Move 1 — structure: interaction shape + prototypes as structural anchors.
An interaction occurs when the characteristics of an item match the taste of
a user, and matrix factorization models mirror that shape when they compute a
recommendation score as the product of a user embedding and an item
embedding. ProtoMF keeps this bilinear form but routes it through prototypes:
users and items are represented by their similarity to a small set of learned
prototype vectors, so the prototypes act as the structural anchors of the
embedding space, and every score is composed from positions relative to them.

% Move 2 — problem + ProtoMF's post-hoc answer vs model-intrinsic.
These anchors, like everything else in the model, are learned from the
collaborative signal, often exclusively and at massive scale. The
collaborative signal records who interacted with what, and nothing about what
any of those things are, so no dimension of a representation learned from it
carries a meaning a human could read off. ProtoMF answers this after the
fact: a trained prototype is interpreted by inspecting its most similar items
and naming it accordingly~\cite{melchiorre2022protomf}. % TODO bib entry
That is a post-hoc reading of an opaque object, and it stands one step short
of the model-intrinsic interpretability introduced in
Section~\ref{sec:bg-iml-intrinsic}: the name is attached to the anchor, it
was never part of it.

% Move 3 — grounding, defined; forced, not offered; vocabulary precedent.
We close that step by grounding the prototypes. Grounding, as used in this
thesis, means constraining the latent factors to a chosen interpretable
vocabulary: the model may only build its internal representations from
ingredients that carry meaning on their own. Merely offering attribute
information alongside the collaborative signal would leave the model free to
ignore it, and would add explanations only as decoration. Constraining the
internal vocabulary itself has precedent outside recommendation, where
restricting a model to expert-defined ingredients produced interpretable
internals at no loss of accuracy~\cite{niknazarXXXXsleep}. % TODO bib entry

% Move 4 — payoff: self-explanatory anchors -> inherently explainable score.
A prototype grounded this way explains itself: its meaning is its
composition, not a name assigned afterwards. And because the prototypes
anchor the space, the property propagates to the score. A recommendation
becomes a statement in the two ingredients of the interaction: the
characteristics an item offers, and the taste a user brings.

% Move 5 — hand-off into the asymmetry.
Grounding, however, needs its vocabulary to come from somewhere. What the two
ingredients offer differs sharply, and this asymmetry shapes the entire
lineage of models in this thesis.

% --- PART TWO: the asymmetry (Matteo's text, agreed repairs applied) ---

% Item side — the packed sentence split (V2-4), coinages lowercased.
For the item ingredient, the vocabulary is already written down. A user
approaches an item with no prior external knowledge; to them, an item simply
is the characteristics shown about it. To the extent that these
characteristics enter the model as interpretable features, they provide the
appropriate vocabulary to explain item prototypes as groups of items sharing
attributes.

% User side — his sentences; unlabeled-vector verb fixed; hedge plainened.
Explaining user prototypes as groups of shared taste is not as
straightforward. Users do not state their taste, and in most cases would not
be able to articulate it if asked. Instead, users reveal their taste through
their interactions. A CF user embedding captures exactly this, but as an
unlabeled vector, with no vocabulary attached that could describe what was
learned and ground prototypes. User attributes correlate with taste only
loosely, are thin in practice, and connect to privacy concerns when used for
explanations; they therefore enter only as a secondary mechanism (see
Section~\ref{sec:fu-cold-users}). Interaction histories must be the primary
signal, adapted into an interpretable vocabulary.

% Transition — granularity fixed, refs in, callback anchor for fU (V2-7).
% (The old "Let it again be made clear..." block dissolved into moves 3-4.)
The next section therefore constructs the easier item side. The harder user
side follows in Chapter~\ref{ch:fu}, where its grounding routes through the
item vocabulary established here, and Chapter~\ref{ch:fufi} merges the two.
